% ================================================================================
% Section 7: Summary and Conclusion
% ================================================================================
\section{Summary and Conclusion}

\begin{frame}{Summary of Main Results}
    \begin{itemize}
        \item \textbf{Example from Brenner and Foster:} $1+5^a = 7^b+7^c$ solved successfully with modulo trick
        \item \textbf{Formal definitions:} First rigorous characterization of "solvable by modulo trick"
        \item \textbf{Lemma framework:} Any sequence of moduli $\equiv$ single modulus (LCM)
        \item \textbf{Theorem 1:} $1+(pq)^a = p^b+q^c$ is non-solvable for coprime $p,q$
        \item \textbf{Theorem 2:} $a^x-a^y = b^z-b^w$ is non-solvable for any positive $a,b$
        \item \textbf{Corollaries:} Validated Brenner and Foster's observations (8.033, 8.037)
    \end{itemize}
\end{frame}

\begin{frame}{Conclusion}
\begin{itemize}
    \item[] \begin{block}{What we accomplished}
        \begin{itemize}
            \item Showed a concrete example of successful modulo trick application from Brenner and Foster
            \item Bridged the gap between heuristic observation and rigorous proof
            \item Provided mathematical foundation for why some equations resist elementary methods
            \item Created diagnostic tool for researchers
        \end{itemize}
    \end{block}

    \item[] \begin{block}{Practical recommendation}
        If an equation has the form $1+(pq)^a = p^b+q^c$ or $a^x-a^y = b^z-b^w$:
        \begin{itemize}
            \item Do \textbf{not} waste time with the modulo trick
            \item Proceed directly to advanced methods (Baker's theory, linear forms in logarithms)
        \end{itemize}
    \end{block}
\end{itemize}
\end{frame}

\begin{frame}{Future Work}
    \begin{itemize}
        \item Find non-trivial equations that \textbf{are} solvable with modulo trick
        \item Develop complete characterization of solvability conditions
        \item Apply framework to other forms of exponential Diophantine equations
        \item Investigate connections to Baker's theory and linear forms in logarithms
    \end{itemize}

    \vspace{0.5cm}
    \textbf{Challenges encountered:}
    \begin{itemize}
        \item Finding general conditions for solvability proved difficult
        \item May require deeper mathematical knowledge than expected
    \end{itemize}
\end{frame}

\begin{frame}{Key References}
    \begin{thebibliography}{9}
        \bibitem{eDe} J. L. Brenner and L. L. Foster, "Exponential Diophantine equations," \textit{Pacific J. Math.}, vol. 101, no. 2, pp. 263-301, 1982.

        \bibitem{wang} L. Wang and R. Tijdeman, "Exponential Diophantine equations with four terms," \textit{Indag. Math.}, vol. 51, pp. 87-96, 1988.

        \bibitem{DEZE199247} M. Deze and R. Tijdeman, "Exponential Diophantine equations with four terms," \textit{Indagationes Mathematicae}, vol. 3, no. 1, pp. 47-57, 1992.
    \end{thebibliography}
\end{frame}

\begin{frame}{Acknowledgments}
    \Large
    \centering
    \begin{itemize}
        \item Assoc. Prof. Dr. Kijti Rodtes (Advisor)
        \item Science Classrooms in University-Affiliated School Project (SCiUS)
        \item Naresuan University Secondary Demonstration School
        \item Our families for their unwavering support
        \item Our friends at SCiUS for the wonderful environment
    \end{itemize}
\end{frame}

\begin{frame}
    \Huge
    \centering
    Thank you for your time

    \vspace{1cm}
    \Large
    Any question?
\end{frame}
